% Part: first-order-logic
% Chapter: natural-deduction
% Section: provability-propositional

% verification of properties of provability needed for maximally
% consistent sets in the completeness chapter.

\documentclass[../../../include/open-logic-section]{subfiles}

\begin{document}

\iftag{FOL}
{\olfileid[pt]{fol}{ntd}{ppr}}
{\olfileid[pt]{pl}{ntd}{ppr}}

\olsection{\usetoken{S}{derivability} e os Conectivos Proposicionais}

\begin{explain}
  Estabelecemos que a relação de !!{derivability}~$\Proves$ da
  dedução natural é forte o suficiente para estabelecer alguns fatos
  básicos envolvendo os conectivos proposicionais, tais como $!A \land !B
  \Proves !A$ e $!A, !A \lif !B \Proves !B$ (modus ponens). Esses
  fatos são necessários para a prova do teorema da completude.
\end{explain}

\begin{prop}\ollabel{prop:provability-land}
  \begin{enumerate}
  \item \ollabel{prop:provability-land-left} Tanto $!A \land !B \Proves
    !A$ quanto $!A \land !B \Proves !B$.
  \item \ollabel{prop:provability-land-right} $!A, !B \Proves !A \land !B$.
  \end{enumerate}
\end{prop}

\begin{proof}
  \begin{enumerate}
  \item Podemos derivar ambos:
    \begin{prooftree}
      \AxiomC{$!A \land !B$}
      \RightLabel{\Elim{\land}}
      \UnaryInfC{$!A$}
      \DisplayProof\qquad\bottomAlignProof
      \AxiomC{$!A \land !B$}
      \RightLabel{\Elim{\land}}
      \UnaryInfC{$!B$}
    \end{prooftree}
  \item Podemos derivar:
    \begin{prooftree}
      \AxiomC{$!A$}
      \AxiomC{$!B$}
      \RightLabel{\Intro{\land}}
      \BinaryInfC{$!A \land !B$}
    \end{prooftree}
  \end{enumerate}
\end{proof}


\begin{prop}\ollabel{prop:provability-lor}
  \begin{enumerate}
  \item $!A \lor !B, \lnot !A, \lnot !B$ é inconsistente.
  \item Tanto $!A \Proves !A \lor !B$ quanto $!B \Proves !A \lor !B$.
  \end{enumerate}
\end{prop}

\begin{proof}
  \begin{enumerate}
  \item Considere a seguinte !!{derivation}:
    \begin{prooftree}
      \AxiomC{$!A \lor !B$}
      \AxiomC{$\lnot !A$}
      \AxiomC{$\Discharge{!A}{1}$}
      \RightLabel{\Elim{\lnot}}
      \BinaryInfC{$\lfalse$}
      \AxiomC{$\lnot !B$}
      \AxiomC{$\Discharge{!B}{1}$}
      \RightLabel{\Elim{\lnot}}
      \BinaryInfC{$\lfalse$}
      \DischargeRule{\Elim{\lor}}{1}
      \TrinaryInfC{$\lfalse$}
    \end{prooftree}
    Esta é !!a{derivation} de~$\lfalse$ a partir das hipóteses não
    descarregadas $!A \lor !B$, $\lnot !A$ e $\lnot !B$.
  \item Podemos derivar ambos:
    \begin{prooftree}
      \AxiomC{$!A$}
      \RightLabel{\Intro{\lor}}
      \UnaryInfC{$!A \lor !B$}
      \DisplayProof\qquad\bottomAlignProof
      \AxiomC{$!B$}
      \RightLabel{\Intro{\lor}}
      \UnaryInfC{$!A \lor !B$}
    \end{prooftree}
  \end{enumerate}
\end{proof}

\begin{prop}\ollabel{prop:provability-lif}
  \begin{enumerate}
  \item \ollabel{prop:provability-lif-left}  $!A, !A \lif !B \Proves !B$.
  \item \ollabel{prop:provability-lif-right}
    Tanto $\lnot !A \Proves !A \lif !B$ quanto $!B \Proves !A \lif !B$.
  \end{enumerate}
\end{prop}

\begin{proof}
  \begin{enumerate}
  \item Podemos derivar:
    \begin{prooftree}
      \AxiomC{$!A \lif !B$}
      \AxiomC{$!A$}
      \RightLabel{\Elim{\lif}}
      \BinaryInfC{$!B$}
    \end{prooftree}

  \item Isso é mostrado pelas duas !!{derivation}s seguintes:
    \begin{prooftree}
      \AxiomC{$\lnot !A$}
      \AxiomC{$\Discharge{!A}{1}$}
      \RightLabel{\Elim{\lnot}}
      \BinaryInfC{$\lfalse$}
      \RightLabel{\FalseInt}
      \UnaryInfC{$!B$}
      \DischargeRule{\Intro{\lif}}{1}
      \UnaryInfC{$!A \lif !B$}
      \DisplayProof\qquad\bottomAlignProof
      \AxiomC{$!B$}
      \RightLabel{\Intro{\lif}}
      \UnaryInfC{$!A \lif !B$}
    \end{prooftree}
    Note que $\Intro{\lif}$ pode, mas não precisa, descarregar a
    hipótese~$!A$.
  \end{enumerate}
\end{proof}

\end{document}
